\documentclass[aspectratio=169]{beamer}
% Shared Beamer configuration for MUS2640 lecture decks (included after \documentclass).
\usepackage[utf8]{inputenc}
\usepackage[T1]{fontenc}
\usepackage{lmodern}
\usepackage{graphicx}
\usepackage{booktabs}

\usetheme{Madrid}
\usecolortheme{dolphin}
\usefonttheme{professionalfonts}

\setbeamercovered{transparent}
\setbeamertemplate{navigation symbols}{}
\setbeamertemplate{footline}[frame number]

\newcommand{\coursename}{MUS2640 · Sensing Sound and Music}
\newcommand{\instituteuo}{University of Oslo · Department of Musicology / RITMO}


\title{Week 7 · Harmony and melody}
\subtitle{\coursename}
\institute{\instituteuo}
\date{}

\begin{document}

\begin{frame}
  \titlepage
\end{frame}

\begin{frame}{Aims}
  \begin{itemize}
    \item Relate \textbf{spectral} views (chroma, pitch tracks) to \textbf{symbolic} descriptions (scales, chords).
    \item Introduce expectation / statistical learning as listeners learn grammars.
    \item Flag cultural frames — Western theory as one toolkit among many.
  \end{itemize}
\end{frame}

\begin{frame}{Roadmap}
  \begin{enumerate}
    \item Pitch and pitch class; chroma features (what machines compute).
    \item Harmony as progression of tension/release (functional sketch).
    \item Melody: contour, motifs, phrase structure (listening-led).
    \item Enculturation \& Western-centred labels — critical use.
  \end{enumerate}
\end{frame}

\begin{frame}{Two representations}
  \begin{itemize}
    \item \textbf{Audio/spectral}: partials, tuning, roughness; continuous pitch traces.
    \item \textbf{Symbolic}: scores/MIDI/discrete pitch classes; reduction layers.
  \end{itemize}
\end{frame}

\begin{frame}{Pedagogy note}
  \begin{itemize}
    \item Chapter is tool-rich — assign \textbf{essential} vs \textbf{stretch} sections explicitly.
    \item Prioritise concepts that unlock student listening goals (by cohort).
  \end{itemize}
\end{frame}

\begin{frame}{In-class}
  \begin{itemize}
    \item Short excerpt: annotate harmonic rhythm vs melodic peaks in plain language first.
    \item Then optionally overlay chord symbols if appropriate to repertoire.
  \end{itemize}
\end{frame}

\begin{frame}{Outside class}
  \begin{itemize}
    \item Focused analysis task on a chosen excerpt; reading as per weekly matrix.
  \end{itemize}
\end{frame}

\begin{frame}{Summary}
  \begin{itemize}
    \item Harmony/melody connects psychoacoustics, culture, and computation.
    \item Next: \textbf{The body} — embodiment, gesture, motion capture ideas.
  \end{itemize}
\end{frame}

\end{document}
